\section{Pinned basis, corrections, and scope}
This continuation starts from the integrated repository revision
\begin{center}\small\ttfamily b6825f8faafff4539f93eae9c4ccb9749b32a855\end{center}
The current mixed-trace proof and its corrected receiver composite were read
before the constructions below \citep{Mixed,Composite}. The sharp receiver
proof, its ES1--ES15 derivation and its ES16--ES37 continuation were also read
\citep{Sharp,EScont}. The actual changes matter: the automatic E-to-M
ordered denominators are $(HsJ,pHJr,pHsr)$; the input of a proper trace is
rational, not already an integral receiver input; the returned cofactor is
not necessarily its residual; the full receiver frame has sixteen choices
of signs, rather than eight choices of complete frames. The published
composite explicitly leaves the bare sorted M-origin fibre unclassified.

We complete that fibre, including cross-channel collisions. We then derive
middle-channel receiver asymptotics and two infinite prime families in the
hard class $1\pmod{840}$. They prove optimality of all three inherited inverse
exponents on that narrower middle domain. We do not claim optimality on the
still smaller image of the mixed-source map. In particular, the already
published ES16--ES17 exterior family with $c=29$ and
$p\equiv3361\pmod{24360}$ is not a new result of this note.
Neither conclusion proves that a witness exists for every prescribed prime.

\section{The original source and its corrected return}
Let $p\equiv1\pmod{24}$ be prime. The mixed source consists of
$a=cq$, $c\mid6$, $q>3$ prime, $p/4<a<p/2$, and every positive rational
ordered pair $y,z$ satisfying $4/p=1/a+1/y+1/z$ and $y+z\in\mathbb Z$.
The pair $(c,q)$ is unique. Put $R=4a-p$. Retain an E word and an exterior
swap bit, or an M word with its existing orientation:
\[
 u\mid a^2,\qquad g=\gcd(a,u),\qquad
 (h,r,s)=(g^2/u,u/g,a/g),\quad a=hrs,\quad u=hr^2.
\]
The two rational tails and quotients are
\begin{align}
 E:&\quad (hs\kappa,phr\kappa),&\kappa&=(pr+s)/R,\label{eq:Es}\tag{1E}\\
 M:&\quad (phs\lambda,phr\lambda),&\lambda&=(r+s)/R.\label{eq:Ms}\tag{1M}
\end{align}
Their sum is integral exactly when $R\mid G^2$, for $G=4u+1$ and $G=p+4u$
respectively. Their common denominator is $R/\gcd(R,G)$. The exact gates
require $R\mid G$. These coordinates are the inherited Type I/II framework
\citep[\S2, Propositions 2.2 and 2.6]{ET}.

For completeness, write $N=y+z$. Reduced denominators coincide at $d$,
and $yz=Npa/R$ gives $d^2=R/\gcd(R,N)$. Thus
$Ry-pa,Rz-pa$ are positive integers with product $(pa)^2$.
Their $p$-valuation pairs $(0,2),(1,1),(2,0)$ recover exactly (1E)--(1M),
including the E swap. Primewise gcd normalization gives the unique displayed
$h,r,s$. This proves completeness before the full gate is imposed.

The mixed return $\rho$ in \citet{Mixed} has the following exact form.
Write $R=\delta t^2$ with $\delta$ the positive squarefree part (odd
exponents retained). For an E input, its carrier character forces $rs\mid6$.
Set
\[
 H=\frac{\delta+1}{4rs},\qquad J=\frac{pr+s}{\delta}.
\]
The resulting middle state is $(h_M,r_M,s_M,\lambda_M)=(H,s,J,r)$.
Since $J\ge r+s>s$, its bare sorted image is
\begin{equation}\label{eq:forwardE}
 \rho_E=(p;HsJ,pHsr,pHJr).
\end{equation}
For an M input, its carrier exponent is zero or two. Choose the small word
$u_0=u$ or $a^2/u$ accordingly, so $u_0\mid c^2\mid36$.
At $a'=(p+\delta)/4$ use the full M word $u_0$, then sort the tails.
The remembered complement restores the ordered version before sorting.
The primewise character and unit-square proofs of these returns are those
of \citet[\S\S2--5]{Mixed}; only the formulas just stated are inputs below.

No independent orientation bit is added to an M word: $u$ and $a^2/u$
already encode its two ordered tails. An E word instead has its separately
retained swap bit. This convention prevents artificial factors of two in
all fibre counts.

\section{The full bare sorted fibre}
Let $W=(p;a,Y,Z)$ be any existing sorted middle witness, with
$a<p<Y<Z$ and $2\le Y/p<Z/p$.
Recover its unique primitive orientation by
\[
 b=Y/p,\quad c=Z/p,\quad g=\gcd(b,c),\quad
 r=b/g<s=c/g,\quad h=a/(rs),\quad\lambda=g/h.
\]
Then
\begin{equation}\label{eq:target}
 a=hrs,\quad Y=phr\lambda,\quad Z=phs\lambda,
 \quad R=4a-p=(r+s)/\lambda.
\end{equation}
Define the literal carrier test
\[
 \operatorname{Car}(A)\iff
 \exists c\in\{1,2,3,6\}:c\mid A,\ A/c>3\text{ prime}.
\]
It returns a unique $(c,q)$ when true. Primality is a finite exact test,
not an assumed supply condition.
For an integer $n=\prod_\ell\ell^{e_\ell}$, also write
\[
 K(n)=\prod_\ell\ell^{\lceil e_\ell/2\rceil};
\]
this is the least positive integer with $n\mid K(n)^2$.

Put $\delta_E=4hr\lambda-1$. Define $\mathcal I_E(W)$ to be empty unless
$r\lambda\mid6$ and $\delta_E$ is squarefree. Otherwise it is exactly
\begin{equation}\label{eq:Efibre}
 \left\{t\mid\operatorname{oddpart}(s):
  \delta_Et^2<p,\quad
  \operatorname{Car}\left(\frac{p+\delta_Et^2}{4}\right)\right\}.
\end{equation}
Also set
\[
 \mathcal V(W)=\{hr^2,hs^2\}\cap\operatorname{Div}(36).
\]
For $\nu\in\mathcal V(W)$, define $\mathcal I_{M,\nu}(W)$ to be empty
unless $R$ is squarefree.  The two possible seeds obey
\[
 \frac{p+4hr^2}{R}=4hr\lambda-1,
 \qquad
 \frac{p+4hs^2}{R}=4hs\lambda-1.
\]
Consequently, for every $\nu\in\mathcal V(W)$ the quantity
$A_\nu=(p+4\nu)/R$ is a positive odd integer.  If $R$ is squarefree, use
\begin{equation}\label{eq:Mfibre}
 \left\{t\mid A_\nu:\ Rt^2<p,\quad
  a_t=\frac{p+Rt^2}{4}=c_tq_t,\quad c_t\mid6,
  \ q_t>3\text{ prime},\quad \nu\mid c_t^2\right\}.
\end{equation}
All such $t$ are odd. The final budget is retained explicitly even where it
also follows from the other integer conditions.

\begin{theorem}[Complete bare fibre]\label{thm:fibre}
The complete ordered rational mixed-source fibre of $W$ is the disjoint union
of the following records.
For each $t\in\mathcal I_E(W)$, set
\[
 R_t=\delta_Et^2,\quad a_t=(p+R_t)/4,\quad
 (h_t,r_t,s_t)=\left(\frac{a_t}{r\lambda},\lambda,r\right),
 \quad u_t=h_t\lambda^2.
\]
Take both original E tail orders. For each
$\nu\in\mathcal V(W)$ and $t\in\mathcal I_{M,\nu}(W)$, take the two M words
\[
 R_t=Rt^2,\quad a_t=(p+R_t)/4,\qquad
 u_t=\nu\quad\hbox{or}\quad a_t^2/\nu.
\]
Thus, exactly,
\begin{equation}\label{eq:fibrecount}
 n_W=|\rho^{-1}(W)|
 =2|\mathcal I_E(W)|+
 2\sum_{\nu\in\mathcal V(W)}|\mathcal I_{M,\nu}(W)|.
\end{equation}
 Every original coordinate, prime exponent and ordering of each enumerated
 preimage is recovered by these formulas and gcd normalization. The maps are
 defined on arbitrary existing
sorted M targets; an empty fibre is an allowed, completely decided result.
\end{theorem}
\begin{proof}
In the E return (\ref{eq:forwardE}), the primitive target parameters force
$(r_t,s_t)=(\lambda,r)$, $H=h$, $J=s$ and $\delta=\delta_E$.
 The identity $p\lambda+r=\delta_Es$ follows directly from
 (\ref{eq:target}). For the reconstructed source one has
 \[
  r(4u_t+1)=\delta_E(s+\lambda t^2),\qquad
  \gcd(r,\delta_Et^2)=1.
 \]
 Indeed, $\delta_E\equiv-1\pmod r$ and
 $t\mid s$ with $(r,s)=1$. Multiplication by the unit $r$ therefore shows
 that the squared trace gate at $\delta_Et^2$ is equivalent to
\[
 \delta_Et^2\mid(\delta_Es)^2.
\]
At a prime of the squarefree $\delta_E$, this is
$v_\ell(t)\le v_\ell(s)$; the same inequality holds at every other prime.
Oddness gives exactly $t\mid\operatorname{oddpart}(s)$.
Conversely, that divisibility gives the full squared gate. Since
$4r\lambda\mid24$, $p\equiv1$ and
 $\delta_E\equiv-1\pmod{4r\lambda}$. Moreover
 $\lambda R=r+s$ and $(r,s)=1$ give $(s,r\lambda)=1$; hence
 $t\mid\operatorname{oddpart}(s)$ gives $(t,4r\lambda)=1$. Every unit
 modulo a divisor of $24$ has square one. Thus $t^2$ preserves
 the source shell congruence. The displayed $h_t$ is integral, $u_t\mid a_t^2$,
and (\ref{eq:forwardE}) returns $W$. The carrier test is precisely the extra
restriction from the general automatic-ray source to the mixed source.

In the M return the target residual is the squarefree part of the old
residual. Its two possible original divisor orientations are exactly
$hr^2,hs^2$. The small pre-complement word must be one of these and divide
$c_t^2\mid36$, which gives $\nu$ and its displayed budget.  Put
$N=p+4\nu=RA_\nu$.  Since $R$ is squarefree, the gate
$Rt^2\mid N^2$ is equivalent prime by prime to $t\mid A_\nu$.
For $\nu=hr^2$ or $hs^2$, respectively, $K(\nu)$ divides $hr$ or $hs$;
hence $A_\nu\equiv-1\pmod{4K(\nu)}$.  Therefore $t$ is a unit modulo
$4K(\nu)\mid24$ and $t^2\equiv1\pmod{4K(\nu)}$.  The explicit condition
$\nu\mid c_t^2$ retains the literal carrier budget.

For the complementary word $u_t'=a_t^2/u_t$, first
$\gcd(a_t,R_t)=1$: a common divisor would divide $4a_t-R_t=p$ while
$0<R_t<p$.  Since $u_t\mid a_t^2$, also $\gcd(u_t,R_t)=1$.  Thus
\[
 p+4u_t'\equiv a_tu_t^{-1}(p+4u_t)\pmod {R_t},
\]
and the multiplier is a unit.  The squared M gates for the complementary
words are therefore equivalent.  The returned small word is $\nu$ at $a$,
so its sorted target is $W$.

Every source in either list is a positive rational trace by (1E)--(1M).
Each list has the displayed inverse. Distinct $t$ give distinct old residuals.
The two M words are different because one contains the carrier $q_t^2$ and
the other does not. Different $\nu$ cannot duplicate an M source: its unique
 $q_t$ determines which of $u_t,a_t^2/u_t$ is the small word. E and M sources
are disjoint by the $p$-valuation pair of their original factor coordinates.
The two E orders differ because equal tails would require $s_t=pr_t$, while
$s_t\le a_t<p$. This proves disjointness and completeness.
\end{proof}

The formula is an executable finite inverse. In (\ref{eq:Efibre}) use
$t^2\le(p-1)/\delta_E$; in (\ref{eq:Mfibre}) use
$t^2\le(p-1)/R$. Every prime-factor occurrence in either divisor list is
retained. There is no fixed-progression or subgroup-support replacement.

\subsection{The missing section and channel collisions}
The general-ray section $t=1$ need not be in the mixed source. At
\[
 W=(1009;255,24216,686120),
\]
the full bare fibre consists of the M words $9,11449$ at
$a_t=321=3\cdot107$, $R_t=275$, $t=5$. The squarefree representative
$a=255$ has no mixed carrier. Thus its fibre has two proper inputs and no
$t=1$ member. The E-origin example
\[
 W=(1129;308,3387,1043196)
\]
has only the two E orientations at $a_t=417=3\cdot139$, $t=7$;
 its would-be E-branch section has $a_t=285$, again outside the mixed domain.

An actual hard-prime cross-channel collision is
\[
 W=(6841;1866,20523,12765306).
\]
It has the two E orientations at $a_t=u_t=1713=3\cdot571$ and the two M
words $3,1160652$ at $a_t=1866=6\cdot311$. The four original source records
produce the same bare sorted target, not four new targets.

For fixed $p$, let $\mathcal S_p$ be the ordered mixed-source records and
write $F_W=\rho^{-1}(W)$ for $W$ in the actual image. The undecorated free
abelian pushforward
\[
 \Phi_p:\mathbb Z[\mathcal S_p]\longrightarrow
 \mathbb Z[\operatorname{Im}\rho_p],\qquad e_T\longmapsto e_{\rho(T)},
\]
has kernel basis $e_T-e_{T_W}$ for
$T\in F_W\setminus\{T_W\}$, where $T_W$ is chosen from that explicitly
computed nonempty fibre. Its rank is $\sum_W(n_W-1)$. For the decorated map
$e_{(T,\xi)}\mapsto e_{(W,\xi)}$, with all
$\xi\in\Xi_W$, the rank is instead
$16\sum_W(n_W-1)$. Neither section uses an unavailable $t=1$ source.

Fix now one frame $\xi\in\Xi_W$. With
$O(W,\xi):V_{\mathrm{labels}}\to E_{\mathrm{recv}}$, define
\[
 L_{W,\xi}:\bigoplus_{T\in F_W}V_{\mathrm{labels}}
 \longrightarrow E_{\mathrm{recv}},\qquad
 (v_T)_T\longmapsto O(W,\xi)\sum_Tv_T.
\]
For positive-definite Hermitian forms $G_T$, put
$M_W=\sum_TG_T^{-1}$. The exact quotient Gram operator is
\begin{equation}\label{eq:metricfibre}
 Q_{W,\xi}=O(W,\xi)^{-*}M_W^{-1}O(W,\xi)^{-1}
 =\left[O(W,\xi)M_WO(W,\xi)^*\right]^{-1}.
\end{equation}
The unique minimum-energy preimage of $z$ is
\[
 v_T=G_T^{-1}M_W^{-1}O(W,\xi)^{-1}z.
\]
For $G_T=I$, this becomes
$v_T=n_W^{-1}O(W,\xi)^{-1}z$. This complex vector-space section is generally
neither an integral coefficient section nor a choice of one arithmetic source.
The formulas follow by a Lagrange multiplier or direct completion of the square;
they change no original metric by decree.

\begin{figure}[ht]
\centering
\begin{tikzpicture}[
  >=Latex,
  every node/.style={font=\scriptsize,align=center},
  box/.style={draw,rounded corners,inner sep=4pt,text width=4.25cm},
  target/.style={draw,very thick,ellipse,inner sep=5pt,text width=3.25cm},
  loss/.style={draw,dashed,rounded corners,inner sep=3pt,text width=4.0cm}
]
\node[box] (E) at (-5.0,1.8)
 {E records\\
 $t\in\mathcal I_E(W)$, $R_t=\delta_Et^2$\\
 two original tail orders};
\node[box] (M) at (-5.0,-1.8)
 {M records\\
 $\nu\in\{hr^2,hs^2\}\cap\operatorname{Div}(36)$,
 $t\in\mathcal I_{M,\nu}(W)$\\
 words $\nu$ and $a_t^2/\nu$};
\node[target] (W) at (0,0)
 {$W=(p;a,Y,Z)$\\
 bare sorted middle target\\
 $n_W=2|\mathcal I_E|+2\sum_\nu|\mathcal I_{M,\nu}|$};
\node[box] (O) at (5.0,0)
 {$O(W,\xi):V_{\rm labels}\to E_{\rm recv}$\\
 $\xi\in\Xi_W$, $|\Xi_W|=16$\\
 $(p,O)$ recovers $(W,\xi)$};
\node[loss] (L) at (0,-3.1)
 {forgotten by $\rho$: old shell, $t$, carrier, word choice,
 and sometimes channel; final sorting also forgets tail order};
\draw[->,thick] (E) -- node[above,sloped] {$\rho$} (W);
\draw[->,thick] (M) -- node[below,sloped] {$\rho$} (W);
\draw[->,thick] (W) -- node[above] {choose $\xi$} (O);
\draw[->,dashed] (E.south east) to[bend right=12] (L.west);
\draw[->,dashed] (M.north east) to[bend left=12] (L.west);
\draw[->,dashed] (W) -- (L);
\end{tikzpicture}
\caption{The exact arithmetic fibre and receiver lift.  The E and M records
are disjoint as source records but may meet at one target.  Adding the full
matrix restores $W$ and the selected square-root frame, not the forgotten
arithmetic record.  The fibre theorem and receiver reconstruction formulas
give the exact maps.}
\label{fig:fibre-decomposition}
\end{figure}


\section{A receiver matrix itself retains the target}
Retain the literal roots $\ell=(p,a,Y,Z)$, $S=\sum\ell_j$ and
\[
 H(X)=-S^{-1}\prod_j(X-\ell_j),\quad A=-S^{-1},\quad
 d_j=H'(\ell_j),\quad \xi_j^2d_j=1.
\]
The original receiver
$O:V_{\mathrm{labels}}\longrightarrow E_{\mathrm{recv}}$, with its sixteen
full-frame square-root choices, has
columns
\begin{equation}\label{eq:receiver}
 O_{\cdot j}=\xi_j(1,-id_j,Ad_j^2+2\ell_jd_j,
                  i(7\ell_j^2d_j-13d_j^2))^{\mathsf T}.
\end{equation}
Its invertibility on actual integral witnesses is the current global-sign
and sharp-growth theorem, not an additional assumption \citep{Sharp}.

\begin{proposition}[Fixed-prime target reconstruction]\label{prop:target}
The map $(W,\xi)\mapsto(p,O)$ is injective on the marked target domain.
Its explicit inverse is
\begin{align}
 \xi_j&=O_{0j},& d_j&=iO_{1j}/O_{0j},\label{eq:matrixreturn}\\
 A&=\frac{O_{20}/O_{00}-2pd_0}{d_0^2},&
 \ell_j&=\frac{O_{2j}/O_{0j}-Ad_j^2}{2d_j}.\notag
\end{align}
 The fourth row, sum normalization, square-root equations,
 $d_j=A\prod_{k\ne j}(\ell_j-\ell_k)$ and the arithmetic domain give the
 remaining exact consistency tests.
\end{proposition}
\begin{proof}
All displayed denominators are nonzero. Dividing each receiver column by
its retained first coordinate gives $d_j$ and then the affine expression in
$A,\ell_j$. The known prime is the root with label zero, determining $A$.
The remaining equations solve for each original root. Substitution into
(\ref{eq:receiver}) proves both compositions. No unlabelled norm or singular
spectrum has been used instead of the matrix.
\end{proof}
Thus Theorem~\ref{thm:fibre} also gives the complete arithmetic fibre of the
bare fixed-prime receiver-valued composite for each selected square-root frame.
The matrix recovers the returned sorted target and those four branches, not a
unique choice among its $n_W$ arithmetic preimages. Root permutation and branch
quotienting are not additional unrecorded identifications. This target
reconstruction does not create a missing target.

\section{A fixed-residual middle family and every first correction}
Fix positive integers $R\equiv3\pmod4$ and $u>0$, with $\gcd(R,u)=1$. At the prime
values satisfying the original conditions
\[
 a=(p+R)/4\in\mathbb Z,\qquad u\mid a^2,\qquad R\mid p+4u,
\]
retain the middle witness, sorted for sufficiently large $p$,
\begin{equation}\label{eq:family}
 (a,Y,Z)=\left(a,\frac{p(a+u)}R,\frac{pa(a+u)}{Ru}\right).
\end{equation}
The identity holds at every positive real parameter before these arithmetic
tests; the asymptotic derivation below uses this exact rational family.
For $p>4(R+u)$ the literal order is $a<p<Y<Z$, so the derivative signs
and branch convention below already hold on an explicit half-line.
Define
\[
 \alpha=\frac1{4R},\qquad \beta=\frac1{16Ru},\qquad
 C=\sqrt{\frac3{16R}}.
\]
Take positive real $\xi_a,\xi_Y$, and
$\xi_p=i/\sqrt{-d_p}$, $\xi_Z=i/\sqrt{-d_Z}$, in original label order.
Then $\epsilon=A^2\prod_{i<j}(\ell_j-\ell_i)\prod_j\xi_j=1$.

\begin{theorem}[Middle spectrum with first corrections]\label{thm:spectrum}
For fixed $R,u$, as $p\to\infty$ along this exact family,
\begin{align*}
 s_1&=20\beta^3p^9\left[1+\frac{6R+14u/5}{p}+O_{R,u}(p^{-2})\right],\\
 s_2&=2\alpha^2p^4\left[1+\frac{-R/2+19u/5}{p}+O_{R,u}(p^{-2})\right],\\
 s_3&=\sqrt2 C p^{3/2}\left[1-\frac{41R}{12p}+O_{R,u}(p^{-2})\right],\\
 s_4&=\frac{\sqrt2}{C}p^{-3/2}\left[1+\frac{11R}{12p}+O_{R,u}(p^{-2})\right],
\end{align*}
and
\begin{equation}\label{eq:detexp}
 \det O=80\alpha^2\beta^3p^{13}
 \left[1+\frac{3R+33u/5}{p}+O_{R,u}(p^{-2})\right].
\end{equation}
In particular,
\[
 p^{-3/2}O^{-1}\longrightarrow\frac C2(-i,1,0,0)^{\mathsf T}e_0^*,
\]
\[
 p^{-3/2}\|O^{-1}\|_2\to\sqrt{3/(32R)},\quad
 \|\wedge^2O^{-1}\|_2\to\tfrac12,\quad
 p^4\|\wedge^3O^{-1}\|_2\to4R^2.
\]
Every remainder is for fixed $R,u$. No claim uniform in unbounded $R,u$ is
made, and a value of $p$ used in an analytic check is not thereby declared prime.
\end{theorem}

\subsection{The complete analytic matrix}
Let $z=1/p$ and define
\[
 B(z)=\operatorname{diag}(1,p^{-3},p^{-4},p^{-6})O_p
       \operatorname{diag}(p^{3/2},p^{3/2},1,p^{-3}).
\]
This is an explicitly compensated calculation, not an identification of the
original norms under rescaling. Its inverse scaling is retained below.
Directly from the rational roots and $H'$, the derivative expansions are
\begin{align*}
 d_p&=-C^2p^3[1-10R/(3p)+O(p^{-2})],&
 d_a&= C^2p^3[1-R/(3p)+O(p^{-2})],\\
 d_Y&=\alpha^2p^4[1-3R/p+O(p^{-2})],&
 d_Z&=-\beta^2p^6[1+4R/p+O(p^{-2})].
\end{align*}
For example, the first coefficients of the three differences at the prime
root are $-R/3,-3R+4u,2R+4u$, while the relative first coefficient of $S$
is $2R+8u$. Their sum minus the latter is $-10R/3$.
The same difference calculation gives the other three coefficients.

After factoring these derivative leading terms out of the square roots,
every remaining radical is analytic and equal to one at $z=0$. Every
entry of $B$ then has a nonnegative integral valuation. Thus
$B(z)=B_0+zB_1+O_{R,u}(z^2)$, where
\[
B_0=\begin{pmatrix}
i/C&1/C&0&0\\
-C&-iC&0&0\\
-2iC&C/2&2\alpha^2&0\\
13C^3&-13iC^3&-6i\alpha^3&20\beta^3
\end{pmatrix},
\]
\[
B_1=\begin{pmatrix}
5Ri/(3C)&R/(6C)&0&0\\
5RC/3&iRC/6&-i\alpha&0\\
i(10RC/3-C^3/\beta)&5RC/12-C^3/\beta&(-R+4u)\alpha^2&-3i\beta^2\\
7C-65RC^3&i(13RC^3/2+7C/16)&i(62R+56u)\alpha^3&(120R+56u)\beta^3
\end{pmatrix}.
\]
These coefficients follow from $(1+ez)^{-1/2}=1-ez/2+O(z^2)$ and the four
receiver polynomials in (\ref{eq:receiver}). Every original column, including
its zero leading entries and next coefficient, is shown.
An explicit analytic neighbourhood can be computed from the finitely many
normalized numerator and denominator polynomials of the derivative units:
for each $P(0)=1$ choose $\rho\le[4(1+\sum_{k>0}|[z^k]P|)]^{-1}$.
Then $|P(z)-1|<1/4$ for $|z|\le\rho$. The chosen square-root branches are
analytic there; continuity and $\det B_0\ne0$ give an inverse neighbourhood.
The printed $O_{R,u}$ terms do not assert a uniform numerical neighbourhood.

\subsection{Complete compound calculation and error control}
One has $\det B_0=80\alpha^2\beta^3$. At exterior ranks $1,2,3,4$, the
leading norms and their relative first coefficients are
\[
\begin{array}{c|c|c|c}
 k&\text{power of }p&\text{leading norm}&\text{first coefficient}\\\hline
1&9&20\beta^3&6R+14u/5\\
2&13&40\alpha^2\beta^3&11R/2+33u/5\\
3&29/2&40\sqrt2 C\alpha^2\beta^3&25R/12+33u/5\\
4&13&80\alpha^2\beta^3&3R+33u/5
\end{array}
\]
Here the first rank uses row 3, column $Z$. The second uses rows 2,3 and
columns $Y,Z$. The third uses the complete row 1,2,3 on both triples
$(p,Y,Z)$ and $(a,Y,Z)$; their leading squared norms are equal and their
relative first coefficients are $-R/6+33u/5$ and $13R/3+33u/5$.
Their average gives the displayed correction. The fourth uses
$\operatorname{tr}(B_0^{-1}B_1)=3R+33u/5$.

To check that no other minor changes those coefficients, the row powers of
$O$ are $(0,3,4,6)$ and the column powers are $(-3/2,-3/2,0,3)$.
Add the appropriate powers for each minor, reducing its order by at least
one when its constant $B$-minor vanishes. At rank 1 all other entries are
$O(p^6)$; at rank 2 all other entries are $O(p^{12})$; at rank 3 all other
entries are $O(p^{13})$. The finite constant/linear minor table supplied
in the certificate includes every row and column combination.
The leading entry or leading row gives a lower bound for the operator norm;
the full Frobenius norm gives the matching upper bound. The squares of the
omitted entries contribute relative $O(p^{-2})$ or smaller. Thus the displayed
first coefficients concern the entire exterior operators, not selected minors.
Successive ratios of compound norms prove Theorem~\ref{thm:spectrum}.

For a second direct calculation, put $z_0=B_0^{-1}e_0$ and
$z_1=-B_0^{-1}B_1z_0$. They are exactly
\[
 z_0=(-iC/2,C/2,9R/8,576iRu^3/5)^{\mathsf T},
\]
\[
 z_1=\left(-2iRC/3,-19RC/12,
 \frac{3R(25R-228u)}{80},
 \frac{24iRu^3(-1160R-717u)}{25}\right)^{\mathsf T}.
\]
Their defining equations verify all coordinates. The original inverse is
\begin{equation}\label{eq:invscale}
 O^{-1}=\operatorname{diag}(p^{3/2},p^{3/2},1,p^{-3})
 B(1/p)^{-1}\operatorname{diag}(1,p^{-3},p^{-4},p^{-6}).
\end{equation}
In particular, its four first-column powers are $p^{3/2},p^{3/2},1,p^{-3}$,
not the exterior fixed-$c$ family's $p^{3/2},p^{3/2},1,p^{-6}$.

\section{The middle observation hierarchy is \texorpdfstring{$0,2,3,9$}{0,2,3,9}}
Let $v_p$ be the right singular vector for $s_4(O)$, with its $a$ coordinate
positive. Set
\[
 v_0=(-i,1,0,0)^{\mathsf T}/\sqrt2,\quad
 w_0=(i,1,0,0)^{\mathsf T}/\sqrt2,\quad n_0=C/\sqrt2.
\]
Then
\[
 v_p=v_0-\frac{9R}{4p}w_0+
 \frac{9R}{8n_0}p^{-3/2}e_Y+O_{R,u}(p^{-2}),
\]
and the two tail components have the stronger relative expansions
\begin{align}
 e_Y^*v_p&=\frac{9R}{8n_0}p^{-3/2}
 \left[1+\frac{7R/4-38u/5}{p}+O_{R,u}(p^{-2})\right],\label{eq:tailsv}\\
 e_Z^*v_p&=\frac{576iRu^3}{5n_0}p^{-9/2}
 \left[1+\frac{-35R/4-239u/40}{p}+O_{R,u}(p^{-2})\right].\notag
\end{align}
To prove the last, small component rather than infer it from a larger vector
error, divide (\ref{eq:invscale}) by $p^{3/2}$ and write
\[
 D(z)B(z)^{-1}E(z),\quad
 D=\operatorname{diag}(1,1,z^{3/2},z^{9/2}),\quad
 E=\operatorname{diag}(1,z^3,z^4,z^6).
\]
The right Gram has a simple leading eigenvalue $C^2/2$. Its complementary
block equations give its top eigenvector proportional to
$(1,O(z^3),O(z^4),O(z^6))$. Multiplication by $E$ leaves the coefficient
$e_0+O(z^6)$. After return to label coordinates the vector is therefore
$D(z)[B(z)^{-1}e_0+O(z^6)]$. Its norm is
$n_0[1-11Rz/12+O(z^2)]$. Substitution of $z_0,z_1$ proves all displayed
coefficients, including the relative error in the last coordinate.

Also
\[
 p^3s_4^2=\frac{32R}{3}+\frac{176R^2}{9p}+O_{R,u}(p^{-2}).
\]
For every fixed nonzero scalar row $L$ and $\tau>0$, put
$\mathcal H_p(\tau)=\exp(-\tau p^3O^*O)$. The complete hierarchy is
\[
\begin{array}{c|c|c}
\text{successive original-label tests}&\nu&
 \lim p^\nu L\mathcal H_p(\tau)L^*\,/\,e^{-(32R/3)\tau}\\\hline
Lv_0\ne0&0&|Lv_0|^2\\
Lv_0=0,
Lw_0\ne0&2&(81R^2/16)|Lw_0|^2\\
Lv_0=Lw_0=0,
Le_Y\ne0&3&(27R^3/2)|Le_Y|^2\\
Lv_0=Lw_0=Le_Y=0&9&(3538944R^3u^6/25)|Le_Z|^2
\end{array}
\]
The four directions are a basis, so the tests exhaust every nonzero row.
The other three heat eigenvalues are $O_{R,u}(e^{-c\tau p^6})$ for a
positive constant at fixed $R,u$; their contribution stays negligible after
these polynomial scalings. Convergence is locally uniform for
$\tau\in(0,\infty)$, not at zero time. Fixed sign changes act in the actual
label frame and transform the tests; phases cannot be replaced by moduli.

\begin{figure}[ht]
\centering
\begin{tikzpicture}[
  >=Latex,
  every node/.style={font=\scriptsize,align=center},
  dot/.style={circle,fill=black,inner sep=1.8pt},
  lab/.style={text width=3.0cm}
]
\draw[->,thick] (-5.5,0) -- (5.6,0) node[right] {power of $p$};
\foreach \x/\e in {-4/-3/2,-1/3/2,1.6/4,5/9}{
  \draw (\x,0.12)--(\x,-0.12) node[below=3pt] {$\e$};
}
\node[dot] (s4) at (-4,0) {};
\node[dot] (s3) at (-1,0) {};
\node[dot] (s2) at (1.6,0) {};
\node[dot] (s1) at (5,0) {};
\node[lab,above=7pt of s4] {$s_4\sim(\sqrt2/C)p^{-3/2}$};
\node[lab,above=7pt of s3] {$s_3\sim\sqrt2C\,p^{3/2}$};
\node[lab,above=7pt of s2] {$s_2\sim2\alpha^2p^4$};
\node[lab,above=7pt of s1] {$s_1\sim20\beta^3p^9$};

\draw[->,thick] (-5.5,-2.2) -- (5.6,-2.2) node[right] {observation power $\nu$};
\foreach \x/\e in {-4.8/0,-2.3/2,0.2/3,5/9}{
  \draw (\x,-2.08)--(\x,-2.32) node[below=3pt] {$\e$};
}
\node[dot] (v) at (-4.8,-2.2) {};
\node[dot] (w) at (-2.3,-2.2) {};
\node[dot] (y) at (0.2,-2.2) {};
\node[dot] (z) at (5,-2.2) {};
\node[lab,above=7pt of v] {$Lv_0\ne0$};
\node[lab,above=7pt of w] {$Lv_0=0$, $Lw_0\ne0$};
\node[lab,above=7pt of y] {$Lv_0=Lw_0=0$, $Le_Y\ne0$};
\node[lab,above=7pt of z] {$L|_{\langle v_0,w_0,e_Y\rangle}=0$};
\end{tikzpicture}
\caption{Two different exact exponent systems for the fixed-$(R,u)$ middle
family.  The upper axis records the four singular-value scales.  The lower
axis records the exhaustive rescalings of
$L\exp(-\tau p^3O^*O)L^*$ for fixed $\tau>0$; its last direction is $e_Z$.
The first corrections and coefficients are stated in the middle-spectrum
theorem and the subsequent observation table.}
\label{fig:middle-spectrum}
\end{figure}


For fixed positive original metrics $G,Q$, the orthonormal representative
$\widehat O=Q^{1/2}OG^{-1/2}$ has inverse
$\widehat O^{-1}=G^{1/2}O^{-1}Q^{-1/2}$. Its leading rank-one limit follows from
Theorem~\ref{thm:spectrum} with both factors retained. Writing
$a_0=C(-i,1,0,0)^{\mathsf T}/2$, the limiting small energy is
$[(a_0^*Ga_0)(e_0^*Q^{-1}e_0)]^{-1}$.
A conductor replaces $Q$ by its actual pullback; neither unitary transport
nor uniform moving-metric constants are inferred. The detailed fixed-metric
method is inherited from ES33--ES37 in \citet{EScont}.

\section{All three upper exponents are sharp on hard middle witnesses}
The current uniform theorem gives, with $\gamma=80/6279^4$, the M bounds
\[
 \|O^{-1}\|_2<\frac{224}{\gamma}p\sqrt w,
 \quad \|\wedge^2O^{-1}\|_2<\frac{180}{\gamma},
 \quad \|\wedge^3O^{-1}\|_2<\frac{41}{4\gamma p^2},
 \qquad w=p/(4a-p)\le p/3
\]
\citep[Theorem 1.1 and middle cofactor proof]{Sharp}.
We now restrict optimality to actual middle states at $p\equiv1\pmod{840}$.

\subsection{First hard progression}
Choose $R=11,u=9$ and primes
\[
 p\equiv8401\pmod{9240}.
\]
They satisfy $p\equiv1\pmod{840}$, $p+36\equiv0\pmod{11}$ and
$3\mid a=(p+11)/4$, so every member has the original M state
(\ref{eq:family}). The residue is coprime to $9240$, hence Dirichlet's theorem
supplies infinitely many such primes
\citep[conclusion at source lines 1257--1265]{DirichletOriginal}; see also
\citet[Theorem 18.1, p.~1]{Dirichlet}.
Theorem~\ref{thm:spectrum} proves positive limits at inverse scales
$p^{3/2}$ and $1$ for its first and second compounds. Its initial prime sample
$p=26881$ has $a=6723=3^4\cdot83$, $R=11$, $u=9$. It has no mixed input
fibre: $r\lambda=68$ excludes the E list, and $(p+36)/11=2447$ is prime,
so only $t=1$ fits the M cutoff, with an unavailable mixed carrier. This is
a concrete reason not to assign these optimality claims to the narrower image.

\subsection{Second hard progression}
For primes
\[
 p\equiv4201\pmod{9240},\qquad s=(p+1)/11,
\]
use the original M state
\[
 (h,r,s,\lambda)=(3,1,(p+1)/11,1),\qquad
 (a,Y,Z)=(3s,3p,3ps).
\]
The residue is again coprime to the modulus and is $1\pmod{840}$.
The identity is immediate from $p=11s-1=12s-s-1$, and the denominators
are sorted for these primes. Put $\gamma_0=3/11$. The three small root
ratios are $1,3/11,3$ and the large root is
$Z=\gamma_0p^2[1+O(p^{-1})]$.
The derivative leading coefficients at those small roots are
\[
 D_0=-16/11,\qquad D_1=240/121,\qquad D_2=60/11.
\]
Take the same positive-real or positive-imaginary branch convention according
to their signs. In rows $0,1,2$, the three core columns, after the respective
row powers $p^{-1},p,p^2$, have the matrix
\[
 C_*=
 \begin{pmatrix}
 i\sqrt{11}/4&11\sqrt{15}/60&\sqrt{165}/30\\
 -4\sqrt{11}/11&-4i\sqrt{15}/11&-2i\sqrt{165}/11\\
 -8i\sqrt{11}/11&24\sqrt{15}/121&12\sqrt{165}/11
 \end{pmatrix}.
\]
Its last row has squared norm $2968784/14641$, the exterior norm squared
of its last two rows is $73267200/161051$, and
$\det C_*=2896/121$. These constants follow by the displayed three-by-three
products; their Gaussian-rational, square-root-cancelled versions are in the
certificate.

The first prime in this progression is $p=4201$, giving
$(a,Y,Z)=(1146,12603,4814346)$. Its mixed-source fibre has exactly the two
M words $3$ and $437772$; this finite occupancy is not used to infer occupancy
for the progression.

The original $Z$ column has last entry
$20\gamma_0^3p^6[1+O(p^{-1})]$. Its other entries are of orders
$p^{-2},p^2,p^4$, while the three core columns have orders
$p^{-1},p,p^2,p^3$. Thus the complete leading compound constants are
\[
20\gamma_0^3p^6,
\quad20\gamma_0^3\sqrt{2968784/14641}\,p^8,
\quad20\gamma_0^3\sqrt{73267200/161051}\,p^9,
\quad20\gamma_0^3(2896/121)p^8.
\]
For ranks two and three, all the entries involving the leading $Z$ last
coordinate form one row, with respectively all core row-2 entries and all
core rows-1,2 minors. Every other term is of lower power. For rank four,
eliminating the $Z$ column changes the core rows only in lower powers.
The same operator/Frobenius comparison therefore proves
\begin{equation}\label{eq:endlimits}
 p^{-1}\|O^{-1}\|_2\to\sqrt{286200/360371},\quad
 \|\wedge^2O^{-1}\|_2\to\frac{\sqrt{2968784}}{2896},\quad
 p^2\|\wedge^3O^{-1}\|_2\to\frac{121}{2896}.
\end{equation}
All constants are nonzero. Along the two reduced prime progressions, these
limits prove that no smaller power than $p^{3/2}$, no decaying bound replacing
$O(1)$, and no faster power decay than $p^{-2}$ is possible in the inherited
three uniform estimates, already on hard middle witnesses.
This does not assert prime-carrier preimages for these families; their
membership in the narrower image of $\rho$ is determined separately by
Theorem~\ref{thm:fibre}. No unproved simultaneous primality condition is used.

\section{Verification, source boundaries, and the remaining theorem}
The main arithmetic replay enumerates all original first-half divisor words
at primes $p\equiv1\pmod{24}$ through 10000. It computes all sorted M targets
and compares their formula (\ref{eq:fibrecount}) fibres, including empty fibres,
against the full forward mixed source. The separate checker imports neither
it nor any predecessor code; it uses positive factor pairs of $(pa)^2$ for
rational inputs and primitive $h,r,s$ enumeration for integral targets.
The exact receiver jets are independently expanded from their literal
rational formulas in standard-library Gaussian-rational Laurent arithmetic;
all leading-compound positions are retained. For the second hard family, the
checker verifies the entered leading $x,D$ data and resulting constants,
while the written calculation derives those data, the sign and the dominance
from the full $p$-dependent receiver. The parameter-grid checks corroborate the
written analytic derivation, not an extrapolated prime theorem.
\begin{verbatim}
python run_all.py --directory reproduced
python fibres.py --target 1009 255 24216 686120
\end{verbatim}
The executed counts, commands and hashes are recorded in the receipts.
There is no independent human mathematical review or Lean build.

The completed source-fibre calculation repairs no false universal claim:
it fills exactly the explicit M-origin gap left by the corrected composite.
The hard-prime middle asymptotics refine a post-witness result; they do not
reverse receiver invertibility into original source occupancy.
The still-unproved statement remains that every prescribed hard prime has
an original E or M state. Neither a nonempty mixed fibre at every M target,
nor a mixed source at every prime, is asserted. The programme remains at
Turn~7 on that existence obligation.
